% Reorganized from the frozen October 15, 2026 supplement.
\section{Area-constrained temporary tasks}
\label{app:area_lower}

In window $e$, suppose $M_e$ candidates have gaps
$\varepsilon-\sigma_{e,k}d_e$, and a new reference makes their signs
irrelevant after that window. Put $w_e=M_eH_e$, $V=\sum_ew_e$, and let
$p_e$ be the product-prior probability that at least half of batch $e$ becomes
determinate simultaneously before its window ends. Define
\[
 \Phi_\delta(u)=
 \begin{cases}
  \operatorname{kl}(1-\delta-u,\delta),&u<1-2\delta,\\
  0,&u\ge1-2\delta,
 \end{cases}
\]
\[
 g_{m,\delta}(p)=\max\left\{0,
 \frac{(p-2\delta)_+}{2}
       \left[\log\frac m{4\log2}\right]_+-1,
 \Phi_\delta(1-p/2)\right\}.
\]
Each displayed component is convex and nondecreasing in $p$.

\begin{proof}[Proof of the general area-to-regret lower bound]
If window $e$ never half-completes, at every round more than half its
candidates remain unresolved. Its area is at least $w_e/2$. On area success,
$\sum_ew_e1_{S_e^c}\le2a$; outside that event this sum is at most $V$.
Consequently
\[
 \sum_ew_ep_e\ge(1-\delta)V-2a.
\]
Stop batch $e$ at first half-completion or at its endpoint. The likelihood
argument above now intersects the product-posterior event with completion
probability $p_e$, giving intersection probability at least
$(p_e-2\delta)_+$ and expected determinate count at least $M_ep_e/2$.
The absolute-likelihood and coordinate-flip derivations therefore give
\[
 \kappa_e\overline EN_e\ge M_eg_{M_e,\delta}(p_e).
\]
For the absolute-likelihood term, when the bracketed logarithm is positive
its Jensen expression is increasing for subset sizes at least $M_e/2$;
when it is nonpositive the stated term is at most zero. Each observation
counted before stopping has regret at least $\varepsilon-d_e$.
The signs of all batches belong to one product prior, and samples of
different batches are disjoint. Hence
\begin{equation}
\label{eq:allocation_lower}
 \sup_\sigma E_\sigma R_T\ge
 \inf_{\substack{0\le p_e\le1\\
       \sum_ew_ep_e\ge(1-\delta)V-2a}}
 \sum_e\frac{(\varepsilon-d_e)M_e}{\kappa_e}
                  g_{M_e,\delta}(p_e).
\end{equation}
This is an information lower relaxation, not a claim that every feasible
allocation is achieved by a causal policy.
\end{proof}

\paragraph{Homogeneous consequence.}
Take $M_e=M$, $H_e=H$, $d_e=d$, and $a=JMH/8$. On success at least $3J/4$
windows half-complete; thus $\sum_ep_e\ge45J/64$. At least $3J/8$ windows
have $p_e\ge1/2$: otherwise the sum is at most
$3J/8+(5J/8)/2=11J/16<45J/64$.
The Bernoulli means lie in $[1/4,3/4]$, giving $\kappa_e\le32d^2$.
The partial-recovery lemma and $\varepsilon-d\ge\varepsilon/2$ yield
\[
 \sup_\sigma E_\sigma R_T\ge
 \frac{3J}{8}\frac{\varepsilon M}{2\cdot1024\cdot32d^2}
             \log\frac M\delta
 \ge2^{-20}\frac{\varepsilon JM}{d^2}\log\frac M\delta.
\]
